\documentclass[letterpaper,11pt]{article}
\usepackage[top=0.8in, bottom=0.8in, left=0.9in, right=0.9in]{geometry}
\usepackage[hidelinks]{hyperref}
\usepackage{lmodern}
\usepackage{parskip}
\usepackage{microtype}
\setlength{\parskip}{6pt}
\setlength{\parindent}{0pt}
\begin{document}
\begin{flushleft}
\textbf{\large Ajay Venkatesh} \\
Brooklyn, NY \\
+1 516 585 9013 \\
\href{mailto:av3855@nyu.edu}{av3855@nyu.edu}
\end{flushleft}
\vspace{10pt}
May 21, 2026
\vspace{6pt}

Hiring Manager \\
Johnson \& Johnson

\vspace{10pt}

Dear Hiring Manager,

My name is Ajay Venkatesh. I'm finishing my M.S. in Computer Engineering at NYU, with production software experience at PwC in Bengaluru, a summer SDE internship at Amazon, and ongoing on-campus work at NYU's Novel AI Technologies. I'm writing about the Associate role in Digital, Data \& Analytics supporting Procurement at Johnson \& Johnson.

The role's emphasis on building reports, dashboards, and data extracts for Procurement stakeholders, plus performing trend and variance analysis on spend and supplier metrics, maps directly onto the data-engineering work I've been doing. At PwC I spent the past few years writing complex queries, stored procedures, and parallelized batch ingestion over \textbf{SQL Server} that materially cut execution time on operational feeds. I built \textbf{Power BI} dashboards over datasets joined across multiple relational databases for business stakeholders, exactly the procurement-style operational analytics the JD describes.

On ETL and large-scale data preparation: at NYU's Novel AI Technologies I engineered \textbf{ETL pipelines} transforming \textbf{NoSQL} into \textbf{PostgreSQL} with concurrency control and locking mechanisms, processing hundreds of thousands of records, then powered \textbf{AWS QuickSight} dashboards for real-time analytics. My Big Data graduate project built a \textbf{PySpark / HDFS} pipeline over millions of flight records with rigorous EDA, anomaly detection, and feature engineering. Identifying data issues, coordinating with technical partners, and validating outputs against source systems is the kind of careful work that procurement analytics demands.

On stakeholder coordination: I've owned features end-to-end alongside architects, product managers, and business stakeholders at PwC, translating ambiguous business questions into reporting and analytical deliverables for non-technical audiences. Clear written and verbal communication has been load-bearing across consulting-style delivery for global engagements.

Stack-wise: \textbf{SQL}, \textbf{Python}, \textbf{PySpark}, \textbf{Pandas}, \textbf{NumPy}, plus \textbf{PostgreSQL}, \textbf{SQL Server}, \textbf{Snowflake}, and \textbf{Power BI}, \textbf{Tableau}, \textbf{QuickSight} for visualization. Comfortable across data warehousing, ETL pipelines, and the procurement-style operational reporting workflows the role describes.

What pulls me toward Johnson \& Johnson is the operational seriousness. Procurement analytics in healthcare runs on accuracy that translates into real supplier decisions and operational outcomes at global scale, where careful, documented work compounds far beyond any single dashboard.

I'd welcome the chance to discuss further. Thanks for considering my application.

\vspace{12pt}
Sincerely, \\
Ajay Venkatesh

\end{document}